% Review package v1 for the instructors (Dudi, Sarit, Alon): the paper's claims, evidence and open decisions,
% before any long text is written. Build: pdflatex review_v1.tex (twice).
\documentclass[10pt]{article}
\usepackage[a4paper,margin=2cm]{geometry}
\usepackage{graphicx}
\usepackage{booktabs}
\usepackage{enumitem}
\usepackage[hidelinks]{hyperref}
\setlist{nosep,leftmargin=*}
\setlength{\parskip}{4pt}
\setlength{\parindent}{0pt}

\title{Transition History for Single-Base-Station 5G Localization\\[2pt]\large Review package v1: claims, evidence and decisions before writing}
\author{Gilad Battat}
\date{October 2026}

\begin{document}
\maketitle

\section*{Purpose}
This is not the paper. It lists what the paper will claim, the evidence for each claim, and five decisions I would like from you before writing the full text (about 12 pages), so that the draft needs no structural rework.

\section*{Draft abstract}
Single-base-station localization from per-report channel measurements is limited by how little one report carries. We study a different way of using the same measurements: the history of a user's reports, i.e.\ its transitions. In three simulated studies (a $15\times15$ controlled grid, a $25\times25$ regression area and a macro cell with 300 users), three propagation settings, two base-station placements and seven model families, adding history lowers the position error consistently, by 14 to 76\,\% for the best model in each setting; only single-antenna devices at macro range, which carry no angle, stay at the level of a constant guess. Sequence models (GRU, 1D-CNN) keep improving up to 15 steps where tree ensembles on the same windows stop at about three; most of their gain is the information in the window and 13 to 28\,\% is its order, and the gain survives missing and unevenly spaced reports. A classical tracking filter on top of a snapshot model is a strong baseline on the grid. How history helps depends on what one report carries: on small cells it helps most without an angle of arrival, at macro range it works by averaging noisy angle reports. Absolute errors are plausible for a single omnidirectional base station and are context, not the claim.

\section*{Claims and evidence}
\begin{enumerate}
\item \textbf{History helps in every setting and model family} (Table~1; macro: 15--21\,\% in seven families, 20 user splits). Exception: single-antenna devices at macro range stay at the constant-guess level, with or without history.
\item \textbf{Sequence models use long windows; trees on the same windows stop at about three steps} (grid, four sets; $25\times25$, Table~2).
\item \textbf{Most of the sequence-model gain is the information in the window, 13--28\,\% is the order} (grid, fair shuffle test; one paragraph, no table).
\item \textbf{The gain survives incomplete and irregular history} if the model is trained with it (grid; Table~3 columns 2--3).
\item \textbf{Placement decides the value of the angle; history helps most where the angle is weak} (AoA: +19.5\,pp at the grid centre, +1.0\,pp beyond the corner).
\item \textbf{At macro range the gain is the averaging of noisy angle reports} (unordered mean recovers it; elevation carries most of it).
\item \textbf{A tracking filter (HMM on a snapshot model) is a strong classical baseline}: on the grid it is better than the GRU (Table~3), also with a generic motion model; reported openly, with the reason (a neighbour-to-neighbour walk at constant speed suits a motion model).
\end{enumerate}

\section*{Table 1. Every study at a glance}
\begin{center}\small\resizebox{\textwidth}{!}{% generated by make_tables.py; do not edit by hand
\begin{tabular}{lllllll}
\hline
Study & Setting & AoA & Best snapshot (m) & Best with history (m) & Constant guess (m) & Error reduction \\
\hline
15$\times$15 grid, mixed & BS centre & no & 5.21 (RF) & 1.24 (GRU, h = 15) & 10.9 & 76 \% \\
15$\times$15 grid, mixed & BS centre & yes & 1.25 (RF) & 0.60 (GRU, h = 15) & 10.9 & 52 \% \\
15$\times$15 grid, mixed & BS corner & no & 4.93 (RF) & 1.42 (GRU, h = 15) & 10.9 & 71 \% \\
15$\times$15 grid, mixed & BS corner & yes & 4.25 (RF) & 1.30 (GRU, h = 15) & 10.9 & 69 \% \\
25$\times$25 regression & UMi LOS & no & 14.5 (RF) & 12.4 (GRU, h = 15) & 37.3 & 14 \% \\
25$\times$25 regression & UMi LOS & yes & 8.2 (RF) & 6.6 (GRU, h = 15) & 37.3 & 20 \% \\
25$\times$25 regression & uniform NLOS & no & 29.3 (RF) & 16.7 (GRU, h = 15) & 37.3 & 43 \% \\
25$\times$25 regression & uniform NLOS & yes & 16.0 (RF) & 8.0 (GRU, h = 15) & 37.3 & 50 \% \\
25$\times$25 regression & mixed & no & 24.3 (RF) & 16.7 (GRU, h = 15) & 37.3 & 31 \% \\
25$\times$25 regression & mixed & yes & 13.3 (RF) & 8.7 (GRU, h = 15) & 37.3 & 35 \% \\
Macro cell, uniform NLOS & multi-antenna & yes & 18.7 (RF) & 15.4 (XGBoost, h = 5) & 37.1 & 18 \% \\
Macro cell, uniform NLOS & single-antenna & no & 34.7 (XGBoost) & 33.1 (RF, h = 5) & 33.6 & 5 \% \\
\hline
\end{tabular}
}\end{center}
\vspace{-6pt}{\footnotesize Mean position error of the best model without history and with history (any depth tested), against always guessing the mean position.}

\section*{Table 2. $25\times25$ regression: best tree model against the GRU}
\begin{center}\small\resizebox{\textwidth}{!}{% generated by make_tables.py; do not edit by hand
\begin{tabular}{llllll}
\hline
Propagation & AoA & Best tree, any depth (m) & CNN h = 15 (m) & GRU h = 15 (m) & GRU gain over best tree (m) \\
\hline
UMi LOS & no & 13.7 (XGBoost, h = 3) & 14.2 & 12.4 ± 0.2 & 1.3 \\
UMi LOS & yes & 8.2 (RF, h = 0) & 8.0 & 6.6 ± 0.1 & 1.7 \\
uniform NLOS & no & 21.7 (XGBoost, h = 15) & 17.2 & 16.7 ± 0.1 & 5.0 \\
uniform NLOS & yes & 11.3 (XGBoost, h = 15) & 8.9 & 8.0 ± 0.0 & 3.3 \\
mixed & no & 19.1 (XGBoost, h = 15) & 17.4 & 16.7 ± 0.1 & 2.3 \\
mixed & yes & 10.8 (XGBoost, h = 15) & 9.5 & 8.7 ± 0.0 & 2.1 \\
\hline
\end{tabular}
}\end{center}

\section*{Table 3. Grid: learned history against a tracking filter, clean and imperfect history}
\begin{center}\small\resizebox{\textwidth}{!}{% generated by make_tables.py; do not edit by hand
\begin{tabular}{llrrr}
\hline
Model & Use of history & Regular (m) & 30 \% missing (m) & Uneven gaps, mean 2 (m) \\
\hline
RF, current report only & no & 4.93 & 4.90 & 4.93 \\
XGBoost, 10 older reports & learned (window) & 3.45 & 4.07 & 3.93 \\
1D-CNN, 10 older reports & learned (sequence) & 2.49 & 3.38 & 3.01 \\
GRU, 10 older reports & learned (sequence) & 1.59 & 2.06 & 1.99 \\
RF + HMM filter, generic motion (2 m) & motion model & 0.93 & 1.16 & 1.41 \\
RF + HMM filter, learned motion & motion model & 0.71 & 0.98 & 1.27 \\
\hline
\end{tabular}
}\end{center}
\vspace{-6pt}{\footnotesize $15\times15$ grid, mixed propagation, BS beyond the corner, RSS and SINR only. Mean position error in metres.}

\section*{Planned figures}
(1) the deployments at the same scale; (2) the 12 grid settings with and without history and AoA; (3) error against history depth for trees, CNN and GRU; (4) the macro gain by channel group (which part of the report carries the gain).

\section*{Decisions I would like from you}
\begin{enumerate}
\item \textbf{Venue and format} (this fixes the length; I assume about 12 pages).
\item \textbf{Is the framing right?} The contribution is the consistent gain from history, not a method or an accuracy; the tracking filter is a baseline that is stronger than the GRU on the grid.
\item \textbf{Robustness to irregular history (claim 4): in the paper or not?}
\item \textbf{One metric throughout} (mean error in metres; grid exact-cell accuracy only as a secondary number)?
\item \textbf{Why history helps without an angle on small cells but not at macro range}: my explanation is that one report on a small cell already narrows the position to a few cells and a short sequence resolves the rest, while at macro range one report carries too little. Does this convince you?
\end{enumerate}

\end{document}
